% Конспект
\ticket{18}{Свойства объема: усиленная монотонность, конечная полуаддитивность}
\label{ticket:18}

\begin{theorem}
Пусть \(\mu:\mathcal{P}\to\mathbb{R}\) --- объём.
Тогда \(\mu\) удовлетворяет следующим свойствам:
\begin{enumerate}
    \item \textbf{Усиленная монотонность.}

    Пусть \(A_1,\ldots,A_n\in\mathcal{P}\) попарно
    не пересекаются, \(A\in\mathcal{P}\) и
    \[
        \bigsqcup_{i=1}^{n}A_i\subset A.
    \]
    Тогда
    \[
        \sum_{i=1}^{n}\mu(A_i)\le\mu(A).
    \]

    \item \textbf{Конечная полуаддитивность.}

    Пусть \(A_1,\ldots,A_n\in\mathcal{P}\),
    \(A\in\mathcal{P}\) и
    \[
        A\subset\bigcup_{i=1}^{n}A_i.
    \]
    Тогда
    \[
        \mu(A)\le\sum_{i=1}^{n}\mu(A_i).
    \]
\end{enumerate}
\end{theorem}

\begin{proof}
\noindent\textbf{1.}
По свойству полукольца
\[
    A\setminus\left(\bigsqcup_{i=1}^{n}A_i\right)
    =\bigsqcup_{j=1}^{\ell}D_j,
    \qquad D_j\in\mathcal{P}.
\]
Тогда
\[
    A=
    \left(\bigsqcup_{i=1}^{n}A_i\right)
    \sqcup
    \left(\bigsqcup_{j=1}^{\ell}D_j\right).
\]
Следовательно,
\[
    \mu(A)
    =\sum_{i=1}^{n}\mu(A_i)
     +\sum_{j=1}^{\ell}\mu(D_j)
    \ge\sum_{i=1}^{n}\mu(A_i).
\]

\medskip
\noindent\textbf{2.}
Пусть
\[
    A\subset\bigcup_{i=1}^{n}A_i.
\]
Положим
\[
    B_i=A_i\cap A\in\mathcal{P}.
\]
Тогда
\[
    A=\bigcup_{i=1}^{n}B_i.
\]
Положим
\[
    C_1=B_1,\qquad
    C_2=B_2\setminus B_1,\qquad
    \ldots,\qquad
    C_k=B_k\setminus
    \left(\bigcup_{j=1}^{k-1}B_j\right),\qquad\ldots
\]
Множества \(C_k\) попарно не пересекаются,
но не обязательно \(C_k\in\mathcal{P}\). При этом
\[
    A=\bigsqcup_{k=1}^{n}C_k.
\]
Имеем конечные разложения
\[
    C_k
    =B_k\setminus\left(\bigcup_{j=1}^{k-1}B_j\right)
    =\bigsqcup_j D_{kj},
    \qquad D_{kj}\in\mathcal{P}.
\]
Тогда
\[
    A=\bigsqcup_{k=1}^{n}C_k
     =\bigsqcup_{k,j}D_{kj},
\]
\[
    \mu(A)
    =\sum_{k,j}\mu(D_{kj})
    =\sum_k\sum_j\mu(D_{kj})
    \le\sum_k\mu(A_k),
\]
\[
    \bigsqcup_j D_{kj}
    =C_k\subset B_k\subset A_k.
\]
\end{proof}